\honest{Lawful Uncertainty}{Eliezer Yudkowsky}{2008}

I quote Robyn Dawes on an old kind of experiment. Seventy per cent of the cards are blue, in random order. Always predicting blue is right 70\% of the time. Subjects instead predict blue about 70\% of the time, which is right only 58\% of the time; paid a nickel per correct guess over a thousand trials, they predicted blue 76\% of the time. \nb{The arithmetic is right. The 76\% could not be checked here.}

This, I say, ``compactly illustrates the most important idea in all of rationality.'' I do not argue that it is the most important.

Subjects ``just keep guessing red.'' Dawes says that despite a thousand trials of feedback they cannot believe they cannot predict. I add that they could have bet blue every time and tested their hunches in private.

``I would suspect,'' I say, that the all-blue strategy ``just didn't occur'' to them. \nb{Experiments by Koehler and James, published in 2010, support this suspicion: most people prefer always betting on the likelier option once both strategies are described to them.}

Three times I say that something ``is a counterintuitive idea'': that the best strategy does not look like the cards, that one should behave lawfully amid randomness, and that one should think lawfully under uncertainty. ``Randomizing your actions takes you further from the target, not closer.'' \nb{Two days later, in ``The Weighted Majority Algorithm,'' Yudkowsky gave the exception: against an opponent who can predict your method, ``it is wise to choose randomly.'' That post is not in this book.}

``And so there are not many rationalists, for most who perceive a chaotic world will try to fight chaos with chaos.'' \nb{The evidence is the card task, and the behaviour there gives way to practice and pay. In a 1961 study of 1000 trials, subjects ended up choosing the likelier option 83\% of the time. In experiments by Shanks and colleagues in 2002, about 70\% of subjects given large payments, regular feedback and long training came to choose it on every trial for at least fifty trials in a row.}

Decision theory, I add, still works against an ``irrational'' opponent. Every bet on red is an expected loss, ``and so too with every departure from the Way in your own thinking.'' I end by asking how many \textsc{Star Trek} episodes, and how many theories of AI, this refutes.
